% !TeX root = main.tex
\chapter{Background}
\label{chap:background}


\section{State of the Art in Speech Translation}
\label{sec:sota}

Speech translation research spans both modular pipelines and direct end-to-end approaches. In the modular tradition, speech is transcribed, translated, and optionally synthesized in separate stages, making it possible to combine specialized ASR, MT, and TTS components that can each be improved independently \cite{bahdanau2014neural,vaswani2017attention,radford2022whisper,wang2020fairseq}. This decomposition remains attractive in practice because it exposes explicit control points for error handling, observability, fallback behavior, and selective optimization.

More recent work on direct speech-to-speech translation demonstrates that tightly integrated models can reduce intermediate representation overhead and potentially improve naturalness in some settings \cite{inaguma2023unity,sarim2024direct,zhang2024direct}. Even so, direct approaches do not eliminate the architectural questions relevant to this thesis. A production-oriented platform still needs to orchestrate ingress, buffering, session context, failure handling, artifact persistence, and evaluation evidence. In other words, model architecture and system architecture solve different layers of the same problem.

Industrial offerings confirm the practical importance of the problem space but do not by themselves resolve the research gap. Cloud platforms from Google, Microsoft, and Amazon provide mature translation services, yet they package architectural decisions behind proprietary boundaries \cite{google_cloud_translate,microsoft_translator,aws_transcribe,aws_translate,aws_polly}. Recent applied studies show that real-time translation is increasingly feasible, but they also underscore the continued need to reason about the whole pipeline rather than isolated models \cite{oinam2025ai,kumar2025realtime}. This thesis therefore positions itself between component-level research and black-box managed services by focusing on the architecture needed to make a transparent, distributed speech translation system workable.

\section{Microservices and ML Maintainability}
\label{sec:microservices_ml}

Microservices are relevant to this thesis because they offer a way to separate concerns that scale differently and fail differently. Newman describes microservices as independently deployable services organized around business or domain capabilities, while Richardson emphasizes explicit boundaries, contracts, and asynchronous interaction patterns as mechanisms for preserving autonomy over time \cite{newman2019monolith,richardson2018microservices}. For a speech translation platform, this matters because audio ingress, speech recognition, machine translation, speech synthesis, orchestration, and persistence rarely share the same performance profile or operational lifecycle.

Maintainability becomes even more important when AI components are involved. Engineering AI systems requires attention not only to code modularity, but also to data movement, model packaging, runtime dependencies, reproducibility, and observability across model-serving boundaries \cite{bass2023engineeringai}. Xu et al. further highlight that scalability and maintainability challenges in machine learning systems often arise from the interaction between model pipelines and their surrounding operational infrastructure, not from the model itself \cite{xu2025scalability}. A microservice architecture does not automatically solve these problems, but it provides a structure in which concerns such as independent scaling, versioned interfaces, and controlled failure domains can be made explicit and studied.

The trade-off is that decomposition introduces coordination overhead. A microservice-based AI pipeline gains modularity and clearer ownership, but it also introduces network hops, serialization costs, and observability requirements that would be less visible in a monolith. This tension between modular maintainability and latency overhead is one of the central architectural trade-offs examined in the thesis.

\section{Event-Driven Architecture (EDA)}
\label{sec:eda}

Event-driven architecture is the coordinating paradigm adopted in this thesis because it decouples producers and consumers while allowing the pipeline to remain observable and extensible. In data-intensive systems, logs and streams provide a durable coordination mechanism that can support replay, asynchronous processing, and bounded service autonomy \cite{kleppmann2017designing}. Enterprise integration patterns further describe how messages, envelopes, and broker-mediated routing can be used to connect heterogeneous components without forcing direct synchronous dependencies \cite{hohpe2004enterprise}.

For speech translation, an event-driven design is attractive because the output of one stage naturally becomes the input of the next. Speech recognition emits text events, translation consumes and enriches those events, and synthesis produces output artifacts for return to the client. This pattern supports service isolation and flexible scaling, while also making it possible to trace processing across stage boundaries through correlation identifiers and hop-to-hop duration measurements. Brokered event streams and contract-driven event exchange are therefore not incidental implementation details; they are part of the architectural rationale for how a distributed AI pipeline can remain diagnosable under load \cite{kleppmann2017designing,richardson2018microservices}.

The same literature also motivates two specific patterns that matter later in the thesis. First, the \emph{claim-check} pattern allows large artifacts such as audio files to be stored out of band while lightweight references move through the event stream \cite{hohpe2004enterprise}. In a secure system, those references must be bounded and short-lived rather than permanent bearer credentials. Second, schema governance is needed so that event producers and consumers can evolve without breaking the pipeline. In the context of this thesis, schema evolution is treated as a high-level contract-governance concern: event changes must preserve interoperability and evidence traceability, not merely serialization success.

\section{AI Pipeline Components}
\label{sec:ai_pipeline_components}

The platform examined in this thesis composes several AI-facing components into one end-to-end workflow. Automatic speech recognition converts audio into text, machine translation transfers meaning across languages, and text-to-speech converts translated text into audible output. Voice activity detection and ingress/orchestration logic support these stages by segmenting input, reducing unnecessary compute, and coordinating the conditions under which downstream services should act.

Each component family has its own research and engineering concerns. Whisper-style ASR emphasizes robustness across noisy and multilingual inputs \cite{radford2022whisper}. Transformer-based MT and related sequence models provide the translation backbone used in many modern systems \cite{bahdanau2014neural,vaswani2017attention}. Speech-to-text and direct speech-translation work show how model architectures can be tuned for latency and multilingual operation, but they do not remove the need for surrounding system decisions about batching, warm starts, queue discipline, and failure containment \cite{wang2020fairseq,inaguma2023unity,zhang2024direct}. The architecture must therefore be studied as the environment in which these AI components become a usable platform.

This component view also supports the thesis' maintainability argument. Treating ASR, MT, TTS, VAD, and orchestration as bounded but cooperating services makes it easier to reason about resource isolation, deployment strategy, and selective optimization. It also creates explicit integration points where observability and evidence collection can be attached without requiring the system to be redesigned as a single model-serving endpoint.

\section{Challenges in Distributed Architectures}
\label{sec:distributed_challenges}

Distributed speech translation is constrained by more than raw model accuracy. Tail latency accumulates across network calls, queueing delays, serialization boundaries, and service startup behavior. A system that appears fast in isolated component benchmarks can still fail to feel real-time when exposed to concurrency, warm-up effects, or failure recovery paths. This is why the thesis treats P99 latency, rather than only average response time, as an architectural driver.

Scalability introduces a second challenge. Different pipeline stages scale according to different resource bottlenecks: ASR and TTS may be GPU-bound, orchestration may be coordination-bound, and persistence or broker infrastructure may become throughput-sensitive under burst load. Distributed-systems literature shows that scaling decisions are inseparable from partitioning strategy, observability, and failure semantics \cite{kleppmann2017designing,bass2021softwarearchitectureinpractice,kubernetes2023}. For this reason, the thesis evaluates scalability as a system property of the whole pipeline rather than as a proxy for the throughput of any single component.

Robustness and diagnosability form the third challenge. A multi-stage AI pipeline must continue to behave predictably when one stage slows down, restarts, or emits malformed data. Correlation identifiers, bounded message contracts, and hop-to-hop timing are essential because they allow failures to be localized and explained. At the same time, observability must remain privacy-safe. In this thesis, performance analytics are treated as metadata-only evidence: aggregated percentiles, event types, and duration summaries are admissible in the manuscript, while raw user content and raw request traces are not. That constraint is part of the research design rather than an afterthought.

\section{Quality Assurance: From Testing to Evaluation}
\label{sec:quality_assurance}

The background for this thesis must distinguish between testing a system and evaluating a research claim. Testing establishes whether the platform behaves as intended under controlled scenarios. Evaluation interprets the resulting evidence against the research questions and success criteria. In a distributed AI pipeline, both are necessary: without testing there is no trustworthy basis for evidence, and without evaluation there is no defensible thesis argument.

Reproducibility is therefore central. The project adopts a uv-first validation workflow so that later claims about latency, robustness, and integration behavior can be tied to repeatable execution rather than undocumented local setups. That workflow matters academically because it narrows the gap between implementation detail and thesis evidence. At the same time, only privacy-safe evidence belongs in the manuscript. This is why the evidence posture distinguishes between shareable metadata and non-shareable debug information, and why later empirical claims must be supported with aggregated measurements and sanitized artifacts rather than raw logs or secret-bearing command examples.

This distinction also clarifies the role of the later evaluation chapter. The background chapter can justify why reproducibility, observability, and privacy-safe evidence handling are necessary. It does not yet claim that the platform satisfies the target metrics; that remains an empirical question for later chapters.

\section{Synthesis and Design Implications}
\label{sec:synthesis_design_implications}

The literature reviewed in this chapter supports a clear design direction. Speech translation is no longer a speculative application area. The research gap has shifted from individual model capability toward architectural composition and empirical system behavior. Microservices provide maintainability and bounded ownership. Event-driven coordination supports loose coupling and diagnosability. Distributed-systems theory explains why latency, scalability, and robustness must be treated together rather than independently.

These observations motivate the design choices developed in the next chapter. The thesis therefore carries forward six implications: (1) the pipeline should be decomposed into explicit service stages; (2) event flow should be traceable across stage boundaries; (3) large artifacts should move through claim-check style references rather than bulk messages; (4) contract evolution must preserve interoperability and evidence continuity; (5) evaluation evidence must be reproducible; and (6) manuscript-facing evidence must remain privacy-safe and aggregated. Together, these implications connect the literature review to the architectural response and make the later chapter sequence answerable against the research questions introduced in Chapter~\ref{chap:introduction}.
